\editionPage[\pdfbookmark[0]{Corrigés des défis}{bm-defi-cor-1-parent}]{defi-cor-1}{D01 · Corrigé}{Cent termes, sans cent calculs}{\hyperref[defi-ex-1]{Énoncé p.\,\pageref*{defi-ex-1}}}
\begin{blocEdition}{defi-cor-1}
\editionEtape{Regrouper sans changer l’ordre}On écrit $(-1+2)+(-3+4)+\cdots+(-99+100)$. Chaque parenthèse vaut 1. Les 100 termes forment 50 paires, donc $\res{S=50}$.
\editionEtape{Retirer le dernier terme}S’arrêter à $-99$ signifie enlever le terme $+100$ de la somme précédente : $50-100=\res{-50}$.
\editionEtape{Contrôle}Avec seulement quatre termes, $-1+2-3+4=2$ : on retrouve deux paires. Le raisonnement repose sur la même structure pour les 50 paires, pas sur une supposition à partir du petit exemple.
\end{blocEdition}
\editionPage{defi-cor-2}{D02 · Corrigé}{Retrouver tout le budget}{\hyperref[defi-ex-2]{Énoncé p.\,\pageref*{defi-ex-2}}}
\begin{blocEdition}{defi-cor-2}
\editionEtape{Suivre les restes}Après la première utilisation, il reste $2/3$ du budget. Après la deuxième, il reste $(3/4)(2/3)=1/2$ du budget. La troisième conserve $2/5$ de ce montant, soit $1/5$ du budget initial.
\editionEtape{Revenir au total}Un cinquième vaut 18 euros, donc le budget est $\res{90}$ euros.
\editionEtape{Calculer les parts}Le tiers initial est 30 euros. Il reste 60 ; son quart est 15. Il reste 45 ; les trois cinquièmes valent 27. Contrôle : $30+15+27+18=90$. Les fractions successives ne portaient pas sur le même total.
\end{blocEdition}
\editionPage{defi-cor-3}{D03 · Corrigé}{Jusqu’où peut-on doubler ?}{\hyperref[defi-ex-3]{Énoncé p.\,\pageref*{defi-ex-3}}}
\begin{blocEdition}{defi-cor-3}
\editionEtape{Modéliser}Au départ, il y a 1 élément. Après $n$ étapes, il y en a $2^n$ : une étape correspond à un nouveau facteur 2.
\editionEtape{Encadrer la capacité}$2^{19}=524\,288$ et $2^{20}=1\,048\,576$. Ainsi $2^{19}<1\,000\,000<2^{20}$.
\editionEtape{Conclure}On peut effectuer $\res{19}$ étapes complètes. Une vingtième dépasserait la capacité. On ne compte pas la valeur de départ comme une étape de doublement.
\end{blocEdition}
\editionPage{defi-cor-4}{D04 · Corrigé}{Quatre carreaux pour un grand carré}{\hyperref[defi-ex-4]{Énoncé p.\,\pageref*{defi-ex-4}}}
\begin{blocEdition}{defi-cor-4}
\editionEtape{Du carreau à l’assemblage}Un carreau a pour côté $\sqrt{18}$ cm. Le grand carré a donc pour côté $\res{2\sqrt{18}\ \mathrm{cm}}$ et pour périmètre $\res{8\sqrt{18}\ \mathrm{cm}}$.
\editionEtape{Un second calcul}L’aire totale vaut $4\times18=72$ cm$^2$, donc le côté est aussi $\sqrt{72}$ cm. Ces deux écritures sont égales : $2\sqrt{18}$ est positif et son carré vaut $4\times18=72$.
\editionEtape{Comparer}Le rapport des côtés est 2, celui des aires 4. Doubler toutes les longueurs ne double pas l’aire : il la multiplie par le carré du rapport.\label{fig-defi-cor-4}\figureGeo[corrige]{defi-racine}

\end{blocEdition}
\editionPage{defi-cor-5}{D05 · Corrigé}{Comparer sans choisir de valeur}{\hyperref[defi-ex-5]{Énoncé p.\,\pageref*{defi-ex-5}}}
\begin{blocEdition}{defi-cor-5}
\editionEtape{Écrire les aires}$A=3(x+4)=3x+12$ et $B=2(x+6)=2x+12$, en cm$^2$.
\editionEtape{Comparer en général}$A-B=(3x+12)-(2x+12)=x$. Puisque $x>0$, cette différence est positive : $\res{A>B}$ pour toutes les longueurs admises.
\editionEtape{Étudier l’égalité}L’égalité des aires imposerait $x=0$, exclu par l’énoncé. Tester plusieurs nombres n’aurait pas prouvé la conclusion pour tous les nombres positifs ; la différence littérale le permet.\label{fig-defi-cor-5}\figureGeo[corrige]{defi-litteral}

\end{blocEdition}
\editionPage{defi-cor-6}{D06 · Corrigé}{Compter les pièces}{\hyperref[defi-ex-6]{Énoncé p.\,\pageref*{defi-ex-6}}}
\begin{blocEdition}{defi-cor-6}
\editionEtape{Choisir une inconnue}Si $x$ est le nombre de pièces de 2 euros, il y a $42-x$ pièces de 1 euro. Leur valeur totale vaut $2x+(42-x)$.
\editionEtape{Résoudre}L’équation $2x+42-x=68$ se réduit en $x+42=68$, donc $\res{x=26}$. Il reste $\res{16}$ pièces de 1 euro.
\editionEtape{Vérifier le contexte}Les nombres sont entiers et non négatifs. On compte $26+16=42$ pièces et $2\times26+16=68$ euros. Contrôler seulement la somme d’argent n’aurait pas vérifié le nombre total de pièces.
\end{blocEdition}
\editionPage{defi-cor-7}{D07 · Corrigé}{Annuler une remise}{\hyperref[defi-ex-7]{Énoncé p.\,\pageref*{defi-ex-7}}}
\begin{blocEdition}{defi-cor-7}
\editionEtape{Raisonner sur des parts}Pour 100 euros au départ, le nouveau prix est 80 euros. Il faut rajouter 20 euros, soit $20/80=1/4=\res{25\,\%}$ du nouveau prix.
\editionEtape{Généraliser}Avec un prix initial $P>0$, le nouveau vaut $0{,}8P$. Une hausse de 25\,\% rajoute $0{,}25\times0{,}8P=0{,}2P$. Le prix final est $0{,}8P+0{,}2P=P$.
\editionEtape{Éviter le piège}Une hausse de 20\,\% n’ajouterait que $0{,}16P$ ; elle donnerait $0{,}96P$. Les deux pourcentages successifs ne portent pas sur la même base.
\end{blocEdition}
\editionPage{defi-cor-8}{D08 · Corrigé}{Un remplissage avec fuite}{\hyperref[defi-ex-8]{Énoncé p.\,\pageref*{defi-ex-8}}}
\begin{blocEdition}{defi-cor-8}
\editionEtape{Première période}Le débit net est $18-3=15$ L/min. En 20 min, le bassin gagne 300 L. Il contient alors $75+300=375$ L.
\editionEtape{Après réparation}Il manque $900-375=525$ L, fournis à 18 L/min. La durée restante vaut $525/18=175/6$ min, soit 29 min 10 s.
\editionEtape{Durée totale}On ajoute les 20 premières minutes : $\res{49\text{ min }10\text{ s}}$. Utiliser un débit moyen non justifié sur les deux phases masquerait le changement dû à la réparation.
\end{blocEdition}
\editionPage{defi-cor-9}{D09 · Corrigé}{Un découpage variable}{\hyperref[defi-ex-9]{Énoncé p.\,\pageref*{defi-ex-9}}}
\begin{blocEdition}{defi-cor-9}
\editionEtape{Établir la formule}L’aire initiale est 180 cm$^2$, donc $\res{A=180-4x^2}$.
\editionEtape{Calculer les trois cas}Pour $x=2,3,4$, les aires restantes valent respectivement 164, 144 et $\res{116}$ cm$^2$. Le choix est $x=4$ cm.
\editionEtape{Retrouver directement}Retirer $180-116=64$ cm$^2$ en quatre carrés identiques donne 16 cm$^2$ par carré. Le côté positif vaut $\sqrt{16}=4$ cm. La contrainte $0<x<5$ est respectée.\label{fig-defi-cor-9}\figureGeo[corrige]{coins-retires}

\end{blocEdition}
\editionPage{defi-cor-10}{D10 · Corrigé}{Même moyenne, autre médiane}{\hyperref[defi-ex-10]{Énoncé p.\,\pageref*{defi-ex-10}}}
\begin{blocEdition}{defi-cor-10}
\editionEtape{Moyennes}Les sommes valent 60 et les effectifs 6. Les deux moyennes valent donc $\res{10}$.
\editionEtape{Médianes}Pour A, les deux valeurs centrales sont 4 et 16 : la médiane choisie est 10. Pour B, elles sont 8 et 8 : la médiane vaut 8.
\editionEtape{Interpréter}Une même moyenne n’impose pas une même médiane ni une même répartition. La valeur 28 compense les petites valeurs de B dans la somme. Pour décrire une série, il faut préciser l’indicateur utilisé et ne pas lui faire dire davantage que ce qu’il mesure.
\end{blocEdition}
\editionPage{defi-cor-11}{D11 · Corrigé}{Une cible et une comparaison exacte}{\hyperref[defi-ex-11]{Énoncé p.\,\pageref*{defi-ex-11}}}
\begin{blocEdition}{defi-cor-11}
\editionEtape{Calculer les aires}L’aire totale vaut 64 cm$^2$, l’aire gagnante $4\pi$ cm$^2$. Le modèle uniforme donne $\res{P=4\pi/64=\pi/16}$.
\editionEtape{Comparer sans arrondi}Comme $\pi<3{,}2$, on a $\pi/16<3{,}2/16=0{,}2=1/5$.
\editionEtape{Interpréter}La probabilité est d’environ 19,6\,\%, légèrement inférieure à 20\,\%. Le calcul porte sur des aires, car le point est uniforme sur la surface. Un rapport de rayons et de côtés ne représenterait pas ce modèle.\label{fig-defi-cor-11}\figureGeo[corrige]{defi-cible}

\end{blocEdition}
\editionPage{defi-cor-12}{D12 · Corrigé}{Une droite, trois rôles}{\hyperref[defi-ex-12]{Énoncé p.\,\pageref*{defi-ex-12}}}
\begin{blocEdition}{defi-cor-12}
\editionEtape{Comparer les triangles}$AB=AC$, $BM=CM$ et AM est commun. Les triangles ABM et ACM sont égaux par le cas CCC.
\editionEtape{Au sommet A}Les angles BAM et MAC se correspondent, donc ils sont égaux. AM est la bissectrice de l’angle A.
\editionEtape{Au point M}Les angles AMB et AMC sont égaux et ont pour somme $180^\circ$ ; chacun vaut $90^\circ$. AM est donc une hauteur. Enfin, c’est une médiane par définition puisqu’elle joint A au milieu de BC. Les trois rôles sont prouvés à partir des données.\label{fig-defi-cor-12}\figureGeo[corrige]{defi-triangles}

\end{blocEdition}
\editionPage{defi-cor-13}{D13 · Corrigé}{Le câble et ses attaches}{\hyperref[defi-ex-13]{Énoncé p.\,\pageref*{defi-ex-13}}}
\begin{blocEdition}{defi-cor-13}
\editionEtape{Distance à franchir}Par Pythagore, $d^2=9^2+12^2=225$, donc $d=15$ m.
\editionEtape{Ajouter les réserves}La longueur totale est $15+2\times0{,}75=\res{16{,}5\ \mathrm m}$. La réserve n’est pas un côté du triangle et n’entre pas sous la racine.
\editionEtape{Calculer le coût}$16{,}5\times4{,}80=\res{79{,}20}$ euros. On suppose que la longueur peut être achetée au dixième de mètre sans supplément ; aucun arrondi de lot n’est demandé.\label{fig-defi-cor-13}\figureGeo[corrige]{defi-pyth}

\end{blocEdition}
\editionPage{defi-cor-14}{D14 · Corrigé}{Une bande entre triangles}{\hyperref[defi-ex-14]{Énoncé p.\,\pageref*{defi-ex-14}}}
\begin{blocEdition}{defi-cor-14}
\editionEtape{Les longueurs}Thalès donne $AN=16/4=\res4$ cm. Par Pythagore, $BC=\sqrt{144+256}=20$ cm. Donc $MN=20/4=\res5$ cm.
\editionEtape{Les aires}ABC a pour aire $12\times16/2=96$ cm$^2$ et AMN $3\times4/2=6$ cm$^2$. La partie extérieure a donc une aire de $\res{90\ \mathrm{cm}^2}$.
\editionEtape{Contrôle}Le rapport des longueurs est $1/4$, mais celui des aires est $6/96=1/16$. L’aire du petit triangle n’est pas le quart de la grande.\label{fig-defi-cor-14}\figureGeo[corrige]{defi-thales}

\end{blocEdition}
\editionPage{defi-cor-15}{D15 · Corrigé}{Choisir la bonne inclinaison}{\hyperref[defi-ex-15]{Énoncé p.\,\pageref*{defi-ex-15}}}
\begin{blocEdition}{defi-cor-15}
\editionEtape{Trois rapports}Les cosinus valent $5/10$, $6/10$ et $7/10$.
\editionEtape{Trois angles}Les angles sont respectivement $60^\circ$, environ $53{,}13^\circ$ et environ $45{,}57^\circ$.
\editionEtape{Choisir}Seul le côté horizontal de $\res6$ m convient à l’intervalle demandé. Le schéma rappelle les deux côtés intervenant dans le rapport. Une mesure au rapporteur sur ce dessin de principe ne remplacerait pas les calculs.\label{fig-defi-cor-15}\figureGeo[corrige]{defi-cosinus}

\end{blocEdition}
\editionPage{defi-cor-16}{D16 · Corrigé}{Deux glissements, une construction}{\hyperref[defi-ex-16]{Énoncé p.\,\pageref*{defi-ex-16}}}
\begin{blocEdition}{defi-cor-16}
\editionEtape{Première image}On ajoute $(4;-2)$ : $A'(4;-1)$, $B'(6;-1)$ et $C'(5;1)$.
\editionEtape{Seconde image}On ajoute ensuite $(-1;5)$ : $\res{A''(3;4),\ B''(5;4),\ C''(4;6)}$. Depuis la figure initiale, le glissement unique ajoute $(3;3)$.
\editionEtape{Aire}Chaque translation conserve l’aire et les longueurs. L’image finale a donc la même aire que le triangle initial, malgré son changement de place. Les deux déplacements sont à reproduire avec leurs sens.\label{fig-defi-cor-16}\figureGeo[corrige]{defi-translation}

\end{blocEdition}
\editionPage{defi-cor-17}{D17 · Corrigé}{Un volume sous deux formes}{\hyperref[defi-ex-17]{Énoncé p.\,\pageref*{defi-ex-17}}}
\begin{blocEdition}{defi-cor-17}
\editionEtape{Volumes initiaux}L’aire de base vaut 36 cm$^2$. Le prisme a pour volume $36\times8=\res{288}$ cm$^3$, la pyramide $288/3=\res{96}$ cm$^3$. Il faut donc trois remplissages pour transférer le même volume de liquide.
\editionEtape{Réduction}Avec un rapport de longueurs $1/2$, le facteur de volume est $1/8$. Les nouveaux volumes sont 36 et 12 cm$^3$.
\editionEtape{Comparer}Le rapport 3 entre les volumes est conservé. Il s’agit de volumes de remplissage : on n’affirme pas que trois objets rigides se rangent nécessairement dans le prisme.\label{fig-defi-cor-17}\figureGeo[corrige]{defi-pyramide}

\end{blocEdition}
\editionPage{defi-cor-18}{D18 · Corrigé}{Une somme programmée et prouvée}{\hyperref[defi-ex-18]{Énoncé p.\,\pageref*{defi-ex-18}}}
\begin{blocEdition}{defi-cor-18}
\editionEtape{Suivre les variables}Les dix tours ajoutent les entiers 1 à 10. À la fin, $\res{s=55}$ et $\res{n=11}$ : la dernière augmentation de n a bien lieu après l’ajout de 10.
\editionEtape{Une preuve sans dix additions}Regrouper $1+10$, $2+9$, $3+8$, $4+7$, $5+6$. Les cinq paires valent 11, donc la somme est $5\times11=55$.
\editionEtape{Repérer l’erreur}Si s était remis à 0 à chaque tour, les sommes précédentes seraient perdues. Au dernier tour, on n’aurait plus que l’ajout de 10, donc s vaudrait 10.\label{fig-defi-cor-18}\figureGeo[corrige]{defi-script}

\end{blocEdition}
